\documentclass[11pt]{article}
\usepackage[utf8]{inputenc}
\usepackage[T1]{fontenc}
\usepackage{geometry}
\usepackage{booktabs}
\usepackage{array}
\usepackage{float}
\usepackage{longtable}
\usepackage{hyperref}
\usepackage{listings}
\usepackage{xcolor}
\usepackage{enumitem}
\usepackage{pdflscape}
\geometry{margin=1in}
\hypersetup{colorlinks=true, linkcolor=blue, urlcolor=blue}
\lstset{
  basicstyle=\ttfamily\scriptsize,
  breaklines=true,
  columns=fullflexible,
  frame=single,
  keepspaces=true
}

\newcommand{\BaseCountryCoreArtists}{2,333}
\newcommand{\AdjacentOnlyArtists}{2,493}
\newcommand{\FullArtistUniverse}{4,824}
\newcommand{\CurrentCountryOnlyArtists}{2,539}
\newcommand{\BillboardAddedArtists}{206}
\newcommand{\CurrentFinalChordRows}{44,058}
\newcommand{\CurrentNewSongRows}{36,584}
\newcommand{\UgArtistsTargeted}{2,333}
\newcommand{\UgChordUniqueSongs}{55,335}
\newcommand{\UgChordJsonTotal}{64,829}
\newcommand{\UgChordJsonExisting}{64,818}
\newcommand{\UgChordJsonMissing}{11}
\newcommand{\UgTabUniqueSongs}{5,243}
\newcommand{\UgTabJsonTotal}{7,299}
\newcommand{\UgTabJsonExisting}{7,296}
\newcommand{\UgTabJsonMissing}{3}
\newcommand{\SupplementAddedArtists}{206}
\newcommand{\SupplementTargetSongs}{1,516}
\newcommand{\SupplementExistingRows}{70,496}
\newcommand{\SupplementRows}{13}
\newcommand{\SupplementFinalRows}{70,504}
\newcommand{\SupplementNewDownloads}{8}
\newcommand{\SupplementAttemptedTargetKeys}{1,087}
\newcommand{\MissingBirthYear}{118}
\newcommand{\MissingBirthCity}{178}
\newcommand{\MissingBirthState}{122}
\newcommand{\MissingBirthCountry}{96}
\newcommand{\MissingWikidataQid}{89}
\newcommand{\MissingMusicbrainzMbid}{531}


\title{Sound of Culture\\Report metodologico sulla pipeline country-only, adjacent e Ultimate Guitar}
\author{Nota tecnica di replica cold-start}
\date{\today}

\begin{document}
\maketitle
\tableofcontents
\clearpage

\section{Obiettivo del report}
Questo secondo report non commenta i trend descrittivi del dataset finale, gia documentati nel report esplorativo, ma ricostruisce in modo operativo l'intera pipeline che porta dal perimetro degli artisti country alla costruzione del dataset chord-level finale. L'obiettivo e duplice: da un lato lasciare una traccia leggibile di \emph{che cosa e stato fatto davvero}; dall'altro rendere la pipeline replicabile sia in modalita snapshot-assisted sia in modalita cold-start, cioe rifacendo gli step chiave a partire dagli script sorgente.

Il punto centrale da chiarire e che il progetto ha piu livelli logici:
\begin{enumerate}[leftmargin=2em]
\item un universo artisti, costruito prima a livello di \texttt{country-only}, poi esteso a \texttt{adjacent-only} e all'unione completa;
\item una pipeline di discovery e download su Ultimate Guitar che usa \texttt{artist\_universe\_country\_only.csv} come input;
\item una fase di supplemento Billboard per artisti o canzoni ancora mancanti;
\item un builder finale che parte dal file \texttt{Restricted Final v6}, aggiunge nuove canzoni UG, riattacca i metadati artista, costruisce i flag Billboard e riempie i \texttt{release\_year} mancanti.
\end{enumerate}

Questa architettura spiega anche perche nel repository coesistono sia script di costruzione ``live'' sia wrapper di replica che ripristinano snapshot e cache. La replica completa, in pratica, e una combinazione tra codice, checkpoint intermedi e output consolidati.

\section{Panoramica generale}
\begin{table}[H]
\centering
\caption{Key scripts in the country-only pipeline}
\label{tab:script_map}
\begin{tabular}{lcc}
\toprule
Script & Role & Main outputs \\
\midrule
build\_country\_artists\_dataset.py & Builds the base country-only, adjacent-only and union artist universes & country\_artists\_*.csv/.dta and QC report \\
enrich\_artist\_universe\_missing\_metadata.py & Top-up pass for residual missing birth metadata in artist universes & artist\_universe\_country\_only.csv/.dta plus checkpoints \\
country\_only\_ultimate\_guitar.py & Discovery, planning and raw JSON download from Ultimate Guitar & discovery JSONs, plan JSON, raw\_tabs\_country, raw\_country\_tabs \\
supplement\_billboard\_country\_chords.py & Adds Billboard artists and missing chart songs not covered by broad UG discovery & supplement discovery JSON and run report \\
build\_country\_only\_chords\_final.py & Merges v6 base, new UG chord rows, chart flags and release-year backfill & Sound\_of\_Culture\_Country\_CountryOnly\_Chords\_Final\_2026\_04\_08.csv/.dta \\
run\_full\_replication.py & Cold-start or snapshot-assisted replication wrapper for country artists & replication package with code, directives and datasets \\
run\_country\_songs\_replication.py & Restores or snapshots the upstream country-song deliverables and caches & song replication package \\
run\_country\_merge\_v6\_replication.py & Restores or rebuilds the restricted final v6 bridge used upstream & merge-v6 replication package \\
\bottomrule
\end{tabular}
\end{table}

\begin{table}[H]
\centering
\caption{Current main artifacts produced by the pipeline}
\label{tab:artifact_counts}
\begin{tabular}{lcc}
\toprule
Artifact & Rows / objects & Comment \\
\midrule
Base country-only artist universe & 2,333 & Country-union core before Billboard additions \\
Adjacent-only artist universe & 2,493 & Adjacent-genre expansion kept separate from the core country list \\
Full artist universe & 4,824 & Country core plus adjacent-only union \\
Current artist\_universe\_country\_only.csv & 2,539 & Current country-only file after Billboard additions and direct missing-data top-up \\
Current new chord-song cache & 36,584 & Songs prepared for append after deduplication and raw JSON parsing \\
Current final country-only chord dataset & 44,058 & Final chord-level dataset used in the exploratory report \\
\bottomrule
\end{tabular}
\end{table}


La fotografia corrente degli artefatti chiave e la seguente:
\begin{itemize}[leftmargin=2em]
\item il nucleo storico \texttt{country-only} contiene \BaseCountryCoreArtists\ artisti;
\item l'universo \texttt{adjacent-only} aggiunge \AdjacentOnlyArtists\ artisti non nel nucleo ma rilevanti come perimetro esteso;
\item l'unione completa degli universi artista arriva a \FullArtistUniverse\ record;
\item il file \texttt{artist\_universe\_country\_only.csv} attualmente in uso contiene \CurrentCountryOnlyArtists\ record, cioe il nucleo country piu \BillboardAddedArtists\ artisti aggiunti nel supplemento Billboard;
\item il dataset finale \texttt{Sound\_of\_Culture\_Country\_CountryOnly\_Chords\_Final\_2026\_04\_08.csv} contiene \CurrentFinalChordRows\ osservazioni;
\item il file intermedio \texttt{country\_only\_chords\_new\_songs\_2026\_04\_08.csv} contiene \CurrentNewSongRows\ righe gia pronte per l'append al v6.
\end{itemize}

La relazione tra questi oggetti e importante: l'universo artisti non e solo un dizionario anagrafico, ma e la chiave che decide quali ricerche fare su UG, quali artisti hanno diritto di entrare nella fase song-level e come attaccare al file finale i metadati demografici recuperati nelle fasi precedenti.

\section{Step 1: costruzione degli universi artista}
\subsection{Logica sostantiva}
La pipeline degli artisti nasce in \texttt{build\_country\_artists\_dataset.py}. Lo script raccoglie seed da Hall of Fame, categorie e liste Wikipedia, amplia il perimetro con query Wikidata sui generi country-correlati, poi standardizza e deduplica. La scelta metodologica qui e chiara: prima si costruisce un universo artista con un criterio relativamente conservativo ma ricco di metadati, poi si usa quell'universo come spina dorsale di tutto il resto.

Nel report QC salvato dallo script, il perimetro iniziale e:
\begin{itemize}[leftmargin=2em]
\item \BaseCountryCoreArtists\ artisti nel \texttt{core country-union};
\item \AdjacentOnlyArtists\ artisti \texttt{adjacent-only};
\item \FullArtistUniverse\ artisti nel full artist universe.
\end{itemize}

Il fatto che il file \texttt{artist\_universe\_country\_only.csv} corrente sia piu grande del nucleo originario non e un errore: e l'effetto della fase successiva di supplemento Billboard, che aggiunge artisti utili per coprire canzoni charted ancora assenti.

\subsection{Come vengono riempiti i campi anagrafici gia nel builder base}
Il builder base recupera prima metadati strutturati da Wikidata, poi li completa con Wikipedia e, solo dove serve, con MusicBrainz. La funzione piu importante e \texttt{prepare\_dataset()}, che:
\begin{enumerate}[leftmargin=2em]
\item unisce i seed con i dettagli Wikidata;
\item riempie gli URL Wikipedia mancanti dalle sitelinks;
\item fa enrichment da infobox e lead Wikipedia;
\item costruisce \texttt{birth\_year} da \texttt{birth\_date} o da indizi secondari;
\item risolve \texttt{birth\_city}, \texttt{birth\_state}, \texttt{birth\_country};
\item applica un fallback da categorie Wikipedia se la geografia e ancora incompleta;
\item usa un enrichment MusicBrainz mirato sui casi ancora incompleti;
\item costruisce score di rilevanza country, flag di campione, confidence e note.
\end{enumerate}

\paragraph{Snippet chiave: orchestrazione del builder artista.}
\lstinputlisting[language=Python]{pipeline_generated/snippets/artist_prepare_dataset.py}

\paragraph{Snippet chiave: risoluzione della geografia di nascita.}
\lstinputlisting[language=Python]{pipeline_generated/snippets/artist_resolve_birth_geography.py}

\paragraph{Snippet chiave: enrichment MusicBrainz mirato nel builder base.}
\lstinputlisting[language=Python]{pipeline_generated/snippets/artist_targeted_musicbrainz.py}

La cosa importante da leggere in questi snippet e l'ordine. Il progetto non fa un'imputazione cieca: parte da campi strutturati forti, usa il lead Wikipedia per colmare lacune narrative, poi applica fallback espliciti e annotabili. Questa gerarchia e una scelta di qualita del dato, non solo di convenienza tecnica.

\section{Step 2: recupero residuo dei missing in \texttt{artist\_universe\_country\_only}}
\subsection{Perche esiste un secondo passaggio}
Il builder base dell'universo artista e pensato per la costruzione generale degli output \texttt{country\_artists\_*}. Ma il file operativo usato a valle dalla pipeline UG e \texttt{artist\_universe\_country\_only.csv}. Per questo esiste uno script separato, \texttt{enrich\_artist\_universe\_missing\_metadata.py}, che fa un \emph{top-up} diretto dei missing residui sul file country-only.

Questa distinzione e cruciale per la replica:
\begin{itemize}[leftmargin=2em]
\item \texttt{build\_country\_artists\_dataset.py} costruisce l'universo artisti generale e i campioni principali;
\item \texttt{enrich\_artist\_universe\_missing\_metadata.py --country-only-direct} rifinisce il file country-only effettivamente usato dagli step successivi.
\end{itemize}

\subsection{Ordine operativo del top-up}
La funzione \texttt{enrich\_country\_only\_direct()} documenta esplicitamente la sequenza del recupero residuo:
\begin{enumerate}[leftmargin=2em]
\item parsing e pulizia della geografia gia presente;
\item risoluzione dei QID Wikidata mancanti dai nomi;
\item enrichment mirato da QID Wikidata;
\item enrichment da Wikipedia;
\item risoluzione dei MBID MusicBrainz mancanti;
\item enrichment MusicBrainz mirato;
\item enrichment MusicBrainz generico;
\item recupero di anno/origine da Wikipedia per i gruppi o i casi residui;
\item enrichment Discogs e iTunes/Apple Music;
\item enrichment dagli official websites via schema JSON-LD;
\item applicazione di override manuali web-confirmed;
\item nuova pulizia geografica e inferenza dello stato da coppie citta-paese gia osservate;
\item refresh delle variabili derivate e salvataggio finale del country-only.
\end{enumerate}

\paragraph{Snippet chiave: pipeline completa di top-up sul country-only.}
\lstinputlisting[language=Python]{pipeline_generated/snippets/artist_country_only_direct.py}

\paragraph{Snippet chiave: risoluzione dei QID mancanti dai nomi.}
\lstinputlisting[language=Python]{pipeline_generated/snippets/artist_resolve_missing_qids.py}

\paragraph{Snippet chiave: fallback Discogs/iTunes e official websites.}
\lstinputlisting[language=Python]{pipeline_generated/snippets/artist_discogs_official.py}

\begin{table}[H]
\centering
\caption{Current residual missingness in artist\_universe\_country\_only.csv}
\label{tab:missingness}
\begin{tabular}{lc}
\toprule
Variable & Remaining missing in current country-only universe \\
\midrule
birth\_year & 118 \\
birth\_city & 178 \\
birth\_state & 122 \\
birth\_country & 96 \\
wikidata\_qid & 89 \\
musicbrainz\_mbid & 531 \\
\bottomrule
\end{tabular}
\end{table}


Questa tabella va letta con attenzione. I missing residui nel file country-only non significano necessariamente fallimento della pipeline: una parte dei record aggiunti dal supplemento Billboard riguarda gruppi, entita non statunitensi o casi con evidenza documentale piu debole. La regola pratica e che la pipeline non forza valori quando la fonte non supera una soglia minima di credibilita.

\subsection{Logica dettagliata per variabile}
\paragraph{\texttt{birth\_year}.}
L'anno di nascita viene recuperato, in ordine, da:
\begin{enumerate}[leftmargin=2em]
\item \texttt{birth\_date} strutturato;
\item campi Wikidata mirati;
\item infobox e lead Wikipedia;
\item origine/anni attivi nei gruppi;
\item Discogs;
\item official websites;
\item override manuali web-confirmed.
\end{enumerate}

\paragraph{\texttt{birth\_city} e \texttt{birth\_state}.}
La citta e lo stato seguono una logica di parsing gerarchico:
\begin{enumerate}[leftmargin=2em]
\item parsing di \texttt{birth\_place\_raw};
\item gerarchie amministrative Wikidata e MusicBrainz;
\item categorie Wikipedia e seed location fallback;
\item parsing di campi \texttt{Origin} o schema degli official websites;
\item inferenza dello stato da combinazioni citta-paese gia osservate in record completi;
\item correzioni manuali su casi noti o sospetti.
\end{enumerate}

\paragraph{\texttt{birth\_country}.}
Il paese viene riempito da:
\begin{enumerate}[leftmargin=2em]
\item place hierarchy strutturata;
\item country raw fields;
\item MusicBrainz area/context;
\item Discogs e official website schema;
\item override manuali.
\end{enumerate}

In sintesi, il top-up non e un semplice merge di fonti, ma una pipeline gerarchica che prova prima fonti strutturate e riutilizzabili, poi fallback piu rumorosi, e lascia traccia nei \texttt{notes}.

\section{Step 3: discovery e download da Ultimate Guitar}
\subsection{Perche il download e separato in discovery, plan e download}
Lo script \texttt{country\_only\_ultimate\_guitar.py} e costruito in tre sottocomandi distinti:
\begin{enumerate}[leftmargin=2em]
\item \texttt{discover}: cerca su UG gli artisti del country-only e salva fino a tre versioni migliori per canzone;
\item \texttt{plan}: misura quante discovery rows esistono e quante raw JSON mancano ancora sul disco;
\item \texttt{download}: scarica i payload JSON grezzi, con resume, shard opzionali e gestione di rate limits e casi 451.
\end{enumerate}

Questa separazione e importante per la replica cold-start, perche consente di fermarsi e riprendere senza perdere lo stato gia acquisito.

\paragraph{Snippet chiave: discovery su UG con checkpoint e resume.}
\lstinputlisting[language=Python]{pipeline_generated/snippets/ug_discovery.py}

\paragraph{Snippet chiave: plan, download e CLI.}
\lstinputlisting[language=Python]{pipeline_generated/snippets/ug_plan_download_cli.py}

\begin{table}[H]
\centering
\caption{Saved Ultimate Guitar discovery/download plan snapshot}
\label{tab:ug_plan}
\begin{tabular}{lcc}
\toprule
Layer & Artists / songs / JSON & Coverage snapshot \\
\midrule
UG chords & 2,333 artists & 55,335 songs; 64,829 JSONs total; 11 missing \\
UG tabs & 2,333 artists & 5,243 songs; 7,299 JSONs total; 3 missing \\
\bottomrule
\end{tabular}
\end{table}


La lettura corretta del plan snapshot e questa:
\begin{itemize}[leftmargin=2em]
\item il progetto ha mirato a \UgArtistsTargeted\ artisti nel paese country-only base;
\item per le \emph{chords} risultano \UgChordUniqueSongs\ canzoni uniche e \UgChordJsonTotal\ JSON discovery, di cui solo \UgChordJsonMissing\ mancanti nel plan salvato;
\item per i \emph{tabs} risultano \UgTabUniqueSongs\ canzoni uniche e \UgTabJsonTotal\ JSON discovery, con \UgTabJsonMissing\ mancanti.
\end{itemize}

Questo spiega due fatti sostantivi. Primo, la pipeline salva piu JSON per canzone perche tiene fino a tre versioni top-rated. Secondo, il file finale chord-level non usa ciecamente un solo payload, ma puo fondere i primi tre quando servono informazioni complementari.

\section{Step 4: supplemento Billboard}
\subsection{Perche serve un supplemento}
Anche una discovery larga su UG puo lasciare fuori:
\begin{itemize}[leftmargin=2em]
\item artisti che non erano nel country-only base ma compaiono nelle liste Billboard;
\item canzoni charted di artisti gia presenti ma non ancora intercettate dalla discovery larga.
\end{itemize}

Per questo esiste \texttt{supplement\_billboard\_country\_chords.py}, che lavora su due fronti:
\begin{enumerate}[leftmargin=2em]
\item full discovery per gli artisti aggiunti dal supplemento;
\item targeted search song-by-song per canzoni Billboard ancora mancanti.
\end{enumerate}

\paragraph{Snippet chiave: supplemento Billboard.}
\lstinputlisting[language=Python]{pipeline_generated/snippets/billboard_supplement_run.py}

\begin{table}[H]
\centering
\caption{Saved Billboard supplement run report}
\label{tab:supplement}
\begin{tabular}{lc}
\toprule
Metric & Value \\
\midrule
Added artists in supplement & 206 \\
Target songs in executed run report & 1,516 \\
Existing discovery rows before supplement & 70,496 \\
Rows added by supplement & 13 \\
Discovery rows after merge & 70,504 \\
New raw chord JSON downloads & 8 \\
Attempted target keys saved & 1,087 \\
\bottomrule
\end{tabular}
\end{table}


Il punto metodologico qui e che il supplemento non sostituisce la discovery base: la estende in modo mirato. Il report di run salvato documenta un caso concreto con \SupplementAddedArtists\ artisti aggiunti, \SupplementTargetSongs\ target song processate nel run salvato, \SupplementRows\ discovery rows aggiunte e \SupplementNewDownloads\ nuovi raw JSON scaricati.

Questa fase spiega anche perche il file country-only corrente contiene \BillboardAddedArtists\ artisti in piu rispetto al nucleo originario: il supplemento Billboard ha modificato il perimetro operativo usato a valle.

\section{Step 5: costruzione del dataset finale chord-level}
\subsection{Meccanica del builder finale}
Lo script \texttt{build\_country\_only\_chords\_final.py} parte dal file \texttt{Restricted Final v6} e lo estende. La logica e la seguente:
\begin{enumerate}[leftmargin=2em]
\item legge le discovery rows chord-only e i raw JSON scaricati;
\item costruisce nuove canzoni da appendere al v6;
\item deduplica le nuove canzoni per UG id e per coppia artista-canzone;
\item riattacca i metadati artista dall'universo country-only;
\item crea bridge rows sintetiche per artisti con raw JSON ma senza presenza nel file combinato;
\item applica \texttt{num\_chord\_json\_used};
\item costruisce o carica la cache Billboard year-end;
\item applica i flag chart-level;
\item riempie i \texttt{release\_year} mancanti;
\item salva CSV e DTA con label.
\end{enumerate}

\paragraph{Snippet chiave: costruzione delle nuove righe song-level da raw JSON.}
\lstinputlisting[language=Python]{pipeline_generated/snippets/final_build_new_song_rows.py}

\paragraph{Snippet chiave: cache Billboard e build finale del dataset.}
\lstinputlisting[language=Python]{pipeline_generated/snippets/final_build_dataset.py}

Ci sono due dettagli da sottolineare.

Primo, la variabile \texttt{genre} nel file finale e \emph{song-level} e viene letta dal payload UG (\texttt{advertising.targeting\_song\_genre}); non e una proprieta ereditata automaticamente dall'universo artisti. Questo e il motivo per cui, a livello esplorativo, il dataset largo contiene anche brani classificati come \texttt{rock}, \texttt{pop} o \texttt{folk} pur partendo da artisti country-only.

Secondo, il builder finale usa fino a tre JSON chord per canzone e fa \emph{deep backfill} dei payload: non e una semplice scelta del top-1, ma una procedura che cerca di estrarre il massimo contenuto informativo dai migliori file disponibili.

\section{Step 6: recupero dei \texttt{release\_year} mancanti}
\subsection{Gerarchia di imputazione}
Il recupero dei \texttt{release\_year} e uno dei punti piu delicati dell'intera pipeline. Lo script non usa un solo servizio, ma una gerarchia molto esplicita. L'ordine e:
\begin{enumerate}[leftmargin=2em]
\item override manuali \texttt{WEB\_CONFIRMED\_RELEASE\_YEAR\_OVERRIDES};
\item cache storiche e master cache derivate dal v6 e da cache precedenti;
\item cache dedicate della run country-only finale;
\item lookup artist-level su iTunes;
\item lookup artist-level su MusicBrainz recordings;
\item lookup artist-level su MusicBrainz releases;
\item lookup artist-level su Deezer;
\item lookup song-bundle su iTunes partendo dalle canzoni ancora mancanti;
\item fallback song-level: Discogs, poi MusicBrainz, poi Wikidata, poi Wikipedia.
\end{enumerate}

\begin{table}[H]
\centering
\caption{Release-year caches consulted by the final chord builder}
\label{tab:release_caches}
\begin{tabular}{l}
\toprule
Release-year cache file \\
\midrule
artist\_release\_track\_cache\_country\_only\_2026\_04\_08.json \\
release\_years\_cache.json \\
release\_years\_cache\_country\_only\_chords\_final\_2026\_04\_08.json \\
release\_years\_cache\_discogs.json \\
release\_years\_cache\_fuzzy.json \\
release\_years\_cache\_internet.json \\
release\_years\_cache\_restricted\_final\_v6.json \\
release\_years\_cache\_wiki.json \\
artist\_release\_track\_cache\_country\_only\_2026\_04\_08.json \\
\bottomrule
\end{tabular}
\end{table}


\paragraph{Snippet chiave: enrichment dei release year.}
\lstinputlisting[language=Python]{pipeline_generated/snippets/final_enrich_release_years.py}

Questa parte va letta come una strategia di minimizzazione dei costi e massimizzazione della precisione:
\begin{itemize}[leftmargin=2em]
\item prima si prova a riusare conoscenza gia salvata;
\item poi si interroga per artista, che e piu efficiente quando mancano molte canzoni dello stesso autore;
\item solo alla fine si scende al song-level, che e piu costoso e piu rumoroso;
\item durante il loop lo script salva continuamente cache dedicate, cosi la run e riprendibile.
\end{itemize}

\section{Replica cold-start e replica da snapshot}
\subsection{Cold-start reale}
Per rifare davvero la pipeline da capo, la sequenza concettuale e:
\begin{enumerate}[leftmargin=2em]
\item costruire i dataset artista base con \texttt{build\_country\_artists\_dataset.py};
\item fare top-up del country-only con \texttt{enrich\_artist\_universe\_missing\_metadata.py --country-only-direct};
\item lanciare \texttt{country\_only\_ultimate\_guitar.py discover};
\item controllare copertura con \texttt{country\_only\_ultimate\_guitar.py plan};
\item scaricare i raw JSON con \texttt{country\_only\_ultimate\_guitar.py download};
\item eseguire \texttt{supplement\_billboard\_country\_chords.py};
\item costruire il dataset finale con \texttt{build\_country\_only\_chords\_final.py}.
\end{enumerate}

I comandi pratici, dal root del progetto, sono:
\begin{lstlisting}[language=bash]
python3 execution/step4_country_artists/build_country_artists_dataset.py
python3 execution/step4_country_artists/enrich_artist_universe_missing_metadata.py --country-only-direct

python3 execution/step1_download/country_only_ultimate_guitar.py discover --max-pages 20 --checkpoint-every 10 --workers 3
python3 execution/step1_download/country_only_ultimate_guitar.py plan
python3 execution/step1_download/country_only_ultimate_guitar.py download --kind both --workers 3 --min-delay 0.6 --max-delay 1.6

python3 execution/step1_download/supplement_billboard_country_chords.py --max-pages-artist 20 --max-pages-song 3 --workers 3 --min-delay 0.6 --max-delay 1.6

python3 execution/step2_digitalize/build_country_only_chords_final.py
\end{lstlisting}

\subsection{Replica assistita da snapshot}
Il repository contiene anche wrapper di replica che servono a due funzioni:
\begin{itemize}[leftmargin=2em]
\item ripristinare output consolidati se non si vuole o non si puo rifare una run live;
\item aggiornare e ricreare i replication package con codice, direttive e output correnti.
\end{itemize}

Lo script piu importante in questa famiglia e \texttt{run\_full\_replication.py}, che per la parte country artists decide se:
\begin{enumerate}[leftmargin=2em]
\item usare gli output gia presenti nel root del progetto;
\item ripristinare uno snapshot bundled;
\item forzare un full rebuild se richiesto con \texttt{--full-rebuild}.
\end{enumerate}

\paragraph{Snippet chiave: wrapper di replica full.}
\lstinputlisting[language=Python]{pipeline_generated/snippets/replication_run_full.py}

I wrapper aggiuntivi fanno invece snapshot o restore di componenti specifiche:
\begin{itemize}[leftmargin=2em]
\item \texttt{run\_country\_songs\_replication.py}: output upstream dei song files e release-year caches;
\item \texttt{run\_country\_merge\_v6\_replication.py}: bridge \texttt{Restricted Final v6};
\item \texttt{run\_artist\_universe\_replication.py}: snapshot e restore del solo universo artista.
\end{itemize}

La distinzione metodologica e importante: una replica ``da snapshot'' non e la stessa cosa di una ricostruzione \emph{live}. La prima verifica la portabilita del pacchetto; la seconda verifica la robustezza della pipeline di costruzione.

\section{Come leggere correttamente il rapporto tra country-only, adjacent e dataset finale}
Una possibile fonte di confusione e la presenza simultanea di tre oggetti:
\begin{enumerate}[leftmargin=2em]
\item \texttt{country-only}: perimetro principale usato per la discovery UG e per il dataset finale qui studiato;
\item \texttt{adjacent-only}: perimetro esteso, utile per coverage, confronto e archivio, ma non usato direttamente per il file chord-level country-only;
\item \texttt{country plus adjacent}: unione dei due, utile come universo allargato o come prodotto documentale separato.
\end{enumerate}

In altre parole, le \emph{adjacent} servono a costruire un perimetro esteso degli artisti, ma il dataset \texttt{Sound\_of\_Culture\_Country\_CountryOnly\_Chords\_Final\_2026\_04\_08} viene costruito sull'universo \texttt{artist\_universe\_country\_only.csv}. Il fatto che alcuni brani nel dataset largo non abbiano \texttt{genre == country} non contraddice questo punto: dipende dal fatto che il genere e letto a livello di canzone dal payload UG.

\section{Conclusione operativa}
Se dovessi spiegare la pipeline in una frase sola, direi cosi: il progetto costruisce prima un'identita affidabile degli artisti, poi usa quella identita per scaricare e pulire i contenuti musicali, poi completa i missing piu delicati con una gerarchia di fonti strutturate e fallback documentati, e infine impacchetta sia gli output finali sia le cache necessarie per rendere la replica ripetibile.

Le implicazioni pratiche sono tre:
\begin{itemize}[leftmargin=2em]
\item per rifare il lavoro da zero bisogna pensare la replica come pipeline multi-step, non come singolo script;
\item per ridurre tempi e fragilita conviene conservare checkpoint intermedi e cache, specialmente per \texttt{release\_year};
\item per valutare il dato finale bisogna sempre distinguere tra perimetro degli artisti e classificazione song-level del payload UG.
\end{itemize}

\appendix
\clearpage
\section{Dove trovare il codice completo}
Questo report include soltanto snippet funzionali. Le copie integrali degli script principali sono salvate nella cartella \texttt{pipeline\_code/} del pacchetto del report, mentre gli artefatti di supporto usati per costruire il report sono salvati in \texttt{pipeline\_assets/}.

\end{document}
